\documentclass[UTF8,12pt,a4paper]{ctexart}

\usepackage{amsmath}
\usepackage{booktabs}
\usepackage{geometry}
\usepackage{tabularx}
\usepackage{array}
\usepackage{enumitem}
\usepackage{fancyhdr}
\usepackage{hyperref}
\usepackage{float}

\geometry{left=2.25cm,right=2.25cm,top=2.25cm,bottom=2.25cm}
\setlength{\parindent}{2em}
\setlength{\parskip}{0.16em}
\linespread{1.15}
\hypersetup{hidelinks}
\renewcommand{\arraystretch}{1.22}

\pagestyle{fancy}
\fancyhf{}
\lhead{几何色散/色散与偏振实验操作评分材料}
\rhead{物理实验B}
\cfoot{\thepage}
\setlength{\headheight}{14.5pt}

\newcolumntype{Y}{>{\raggedright\arraybackslash}X}
\newcommand{\scorepoint}[1]{\noindent\textbf{#1}\quad}

\begin{document}

\begin{center}
    {\zihao{2}\bfseries 几何色散/色散与偏振实验操作评分提交材料}
\end{center}

\vspace{0.2em}

\begin{center}
\small
\begin{tabularx}{0.94\textwidth}{@{}>{\bfseries}rY>{\bfseries}rY@{}}
姓名：& 纪文龙 & 学号：& 2024012842 \\
班级：& 行健烽火 4 & 课程：& 物理实验B \\
提交用途：& \multicolumn{3}{l@{}}{操作评分窗口：课堂操作完成情况、原始数据摘录、数据处理与课后总结} \\
\end{tabularx}
\end{center}

\section{提交定位与评分对应}

本材料面向课程平台中的“操作评分窗口”，不是完整实验报告替代件。根据课程说明，剩余实验采用“9+1模式”：实验内容完成度占主要部分，课堂操作表现占 1 分，课后总结可作为对实验理解和操作反思的补充。因此本材料按“可评分证据链”组织：先列明实验完成项，再给出关键操作、原始数据、计算过程、误差分析和课后总结。

\begin{table}[H]
\centering
\caption{本材料与操作评分点的对应关系}
\small
\begin{tabularx}{0.96\textwidth}{p{0.18\textwidth}p{0.18\textwidth}Y}
\toprule
评分关注点 & 本材料位置 & 证明方式 \\
\midrule
实验内容完成度 & 第2、3节 & 五个课堂环节均列出完成内容、操作动作和可核查证据，覆盖最小偏向角、布儒斯特角、虹霓、光谱仪和特斯拉线圈。 \\
课堂操作表现 & 第3、6节 & 说明等高共轴、调焦、找最小偏向角、找反射最暗位置、手动曝光、校准峰选择和高压/激光安全规范。 \\
原始数据与处理 & 第4、5节 & 摘录最小偏向角法和布儒斯特角法的原始角度读数，给出折射率与不确定度计算。 \\
课后总结 & 第7节 & 结合操作难点、误差来源和改进方法总结实验认识，不只停留在现象描述。 \\
\bottomrule
\end{tabularx}
\end{table}

\clearpage
\section{实验内容完成情况总表}

\begin{table}[H]
\centering
\caption{课堂操作完成情况与可评分证据}
\small
\begin{tabularx}{0.96\textwidth}{p{0.17\textwidth}p{0.33\textwidth}Y}
\toprule
实验环节 & 完成内容 & 操作证据/评分点 \\
\midrule
最小偏向角法 & 搭建氦灯、狭缝、凸透镜、转台和望远镜光路；观察氦光谱；测量两条谱线的最小偏向角。 & 能完成等高共轴、调焦、寻找谱线、判断谱线反向运动临界位置、锁紧转台并读出入射/出射方位。 \\
布儒斯特角法 & 用绿光激光照射三棱镜光学面，调节入射角和偏振片，寻找反射光最暗位置。 & 能先确定法线方位，再通过偏振片与转台角度的联调找到最暗点，并用 \(n=\tan\theta_B\) 处理。 \\
虹和霓观察 & 使用玻璃球和亚克力棒观察虹、霓及其偏振特性。 & 能区分一次反射形成虹、二次反射形成霓；能用偏振片验证虹和霓中 S 光占主导。 \\
光谱仪实验 & 启动光谱仪软件，调节选区、曝光和校准峰，观察不同光源谱线。 & 能取消自动曝光、手动调节主峰不过曝；能使用已知氦谱线校准波长并观察其他特征峰。 \\
特斯拉线圈 & 观察电弧放电、无线点亮灯泡/放电管和播放音乐现象。 & 能说明高频高压放电、电磁感应、气体电离和等离子体发声的基本机制，并遵守安全距离。 \\
\bottomrule
\end{tabularx}
\end{table}

\section{关键操作过程记录}

\subsection{最小偏向角法}

\begin{enumerate}[label=\arabic*.,leftmargin=2.2em]
    \item 按“氦灯—狭缝—凸透镜—三棱镜转台—望远镜”的顺序摆放器件，先通过目测和卷尺进行等高共轴粗调。
    \item 将氦灯靠近狭缝，调节凸透镜与狭缝距离，使望远镜中可见清晰、无明显视差的狭缝像。狭缝宽度保持适中，避免谱线过暗。
    \item 放入等边三棱镜后，先用肉眼在出射侧寻找彩色谱线，再将望远镜转到谱线方向，通过十字丝对准目标谱线。
    \item 缓慢转动载有三棱镜的转台，观察谱线在望远镜视场中的运动方向。当谱线即将反向运动时，即为最小偏向角附近。
    \item 锁紧转台，转动望远镜使分划板竖直线与目标谱线重合，读取出射方位角 \(\varphi\)。
    \item 取下三棱镜，使望远镜直接对准狭缝像，读取入射方位角 \(\varphi_0\)。由 \(\delta=\varphi-\varphi_0\) 计算最小偏向角。
\end{enumerate}

\subsection{布儒斯特角法}

\begin{enumerate}[label=\arabic*.,leftmargin=2.2em]
    \item 将三棱镜固定在转台上，用绿光激光照射其光学面，调节到反射光与入射方向重合，以确定法线方位。
    \item 在激光器和三棱镜之间放置偏振片，预设入射角约 \(57^\circ\)，观察白屏上的反射光斑。
    \item 同时微调转台角度和偏振片方位，使反射光斑达到最暗。这里的“最暗”不是绝对消失，而是 P 分量反射最弱的位置。
    \item 读取最暗位置方位，与法线方位作差得到布儒斯特角 \(\theta_B\)，再由 \(n=\tan\theta_B\) 求折射率。
\end{enumerate}

\subsection{虹、霓、光谱仪与特斯拉线圈}

虹和霓部分主要考查几何光学现象识别：玻璃球中一次内反射形成主虹，二次内反射形成副虹，二者颜色排列相反且副虹较暗；转动偏振片可看到亮度变化，说明出射光具有明显偏振特性。光谱仪部分主要考查数字化分光测量的操作规范：需要手动调曝光、选择合适的校准峰，并保持入射光进入光谱仪通光孔。特斯拉线圈部分主要考查高频高压放电现象观察和安全意识：应远离手机、手环、示波器和信号发生器，观察后及时关闭并保持通风。

\section{关键原始数据与处理结果}

\subsection{最小偏向角法测三棱镜折射率}

等边三棱镜顶角取 \(A=60^\circ\)。课堂记录的关键角度读数如下。

\begin{table}[H]
\centering
\caption{最小偏向角法原始角度读数}
\small
\begin{tabular}{ccccc}
\toprule
谱线 & 波长/nm & 出射方位 \(\varphi\) & 入射方位 \(\varphi_0\) & 最小偏向角 \(\delta\) \\
\midrule
紫色 & 447.1 & \(295.4^\circ\) & \(242.0^\circ\) & \(53.4^\circ\) \\
黄色 & 587.6 & \(258.8^\circ\) & \(205.5^\circ\) & \(53.3^\circ\) \\
\bottomrule
\end{tabular}
\end{table}

最小偏向角条件下，三棱镜折射率为
\[
    n=\frac{\sin[(A+\delta)/2]}{\sin(A/2)} .
\]

紫色谱线：
\[
 n_{\text{紫}}=
 \frac{\sin[(60^\circ+53.4^\circ)/2]}{\sin30^\circ}
 =\frac{\sin56.70^\circ}{0.5}
 \approx 1.672 .
\]

黄色谱线：
\[
 n_{\text{黄}}=
 \frac{\sin[(60^\circ+53.3^\circ)/2]}{\sin30^\circ}
 =\frac{\sin56.65^\circ}{0.5}
 \approx 1.671 .
\]

\begin{table}[H]
\centering
\caption{最小偏向角法折射率结果}
\small
\begin{tabular}{cccc}
\toprule
谱线 & 波长/nm & 折射率 \(n\) & 结果判断 \\
\midrule
紫色 & 447.1 & 1.672 & 短波长折射率略大 \\
黄色 & 587.6 & \(1.671\pm0.001\) & 作为不确定度估算对象 \\
\bottomrule
\end{tabular}
\end{table}

\subsection{不确定度估算}

以黄色谱线为例，角度估读精度取 \(0.1^\circ\)。由于 \(\delta=\varphi-\varphi_0\)，若两次角度读数独立，则
\[
    \Delta \delta=\sqrt{(0.1^\circ)^2+(0.1^\circ)^2}
    =0.141^\circ
    =0.00246\ \text{rad}.
\]
由
\[
    \frac{\partial n}{\partial \delta}
    =\frac{\cos[(A+\delta)/2]}{2\sin(A/2)}
\]
可得在 \(\delta=53.3^\circ\) 时
\[
    \Delta n
    \approx
    \frac{\cos56.65^\circ}{2\sin30^\circ}\Delta\delta
    \approx 0.5497\times0.00246
    \approx 0.001 .
\]
因此黄色谱线结果写为
\[
    n_{\text{黄}}=1.671\pm0.001 .
\]
紫色谱线折射率略大于黄色谱线，符合正常色散中短波长光折射率更大的规律；但两条谱线差异与不确定度同量级，因此本组数据更适合作为定性验证，而不宜夸大为高精度定量结论。

\subsection{布儒斯特角法测折射率}

\begin{table}[H]
\centering
\caption{布儒斯特角法原始数据与结果}
\small
\begin{tabular}{ccccc}
\toprule
光源 & 法线方位 & 最暗位置方位 & \(\theta_B\) & \(n=\tan\theta_B\) \\
\midrule
绿光激光 & \(58.6^\circ\) & \(116.0^\circ\) & \(57.4^\circ\) & 1.564 \\
\bottomrule
\end{tabular}
\end{table}

计算过程为
\[
    \theta_B=116.0^\circ-58.6^\circ=57.4^\circ,
\]
\[
    n=\tan57.4^\circ\approx1.564 .
\]
该方法的主要误差来自“最暗位置”的主观判断、偏振片角度不完全准确以及棱镜表面状态。与最小偏向角法相比，布儒斯特角法更依赖观察者对光斑亮度的判断，所以结果可作为对折射率量级和偏振规律的验证，而不是替代最小偏向角法的高精度结果。

\section{现象解释与结果分析}

\scorepoint{色散规律。}
最小偏向角法得到 \(n_{\text{紫}}>n_{\text{黄}}\)，说明短波长紫光相对黄色光偏折更大，这与透明介质的正常色散规律一致。实验中差值较小，原因是两条谱线波长虽有差别，但角度读数精度只有约 \(0.1^\circ\)，折射率差异容易被读数误差淹没。

\scorepoint{偏振判据。}
布儒斯特角实验中，反射光最暗时 P 分量反射最弱，反射光主要为 S 偏振。该现象说明实验不是简单“找一个角度”，而是要同时满足入射角接近布儒斯特角、偏振方向接近 P 偏振、反射光观察屏位置合适等条件。

\scorepoint{虹和霓。}
主虹来源于球形介质内的一次反射，副虹来源于二次反射，因此二者颜色排列相反，副虹亮度更弱。用偏振片观察亮度变化，可以验证虹和霓并非普通漫反射光，而是具有较强的偏振特征。

\scorepoint{光谱仪测量。}
光谱仪测量中，曝光过强会使峰顶饱和、峰位不易判断；曝光过弱则信噪比不足。校准峰若选错，后续所有谱线波长会产生系统偏移。因此操作重点在于手动曝光、准确选峰和记录特征峰，而不是只看软件是否显示曲线。

\scorepoint{特斯拉线圈。}
特斯拉线圈通过高频高压电场使空气或放电管内气体电离，可观察到放电、电磁感应点亮和等离子体发声等现象。该环节虽然不以数值测量为主，但能够体现对电磁振荡、高压放电和安全操作边界的理解。

\section{课堂操作规范与安全注意}

\begin{enumerate}[label=\arabic*.,leftmargin=2.2em]
    \item 光学平台调节应先粗调等高共轴，再细调狭缝、凸透镜和望远镜焦距。狭缝不能调得过窄，否则谱线变暗且难以读数。
    \item 寻找最小偏向角时应缓慢转动转台，重点观察谱线“即将反向运动”的临界位置，锁紧转台后再读数，避免边转边读。
    \item 激光实验中不能直视激光束，也不能让镜面反射光扫过他人视线；氦放电管有高压端，不能触碰电极。
    \item 光学元件只接触边缘，不用手触碰镜面和棱镜光学面；器件不用时放回指定位置，避免划伤和污染。
    \item 光谱仪测量时应手动调曝光，使主峰尖锐但不过曝；校准峰必须选准，否则后续谱线波长会整体偏移。
    \item 特斯拉线圈操作时应远离手机、手环、示波器和信号发生器；手持灯泡不应过近，观察结束后及时关闭并保持通风。
\end{enumerate}

\section{课后总结}

本次实验最能体现操作能力的地方不是公式本身，而是装置调节和判据判断。最小偏向角法看似只需要两个角度，但如果光路没有等高共轴，谱线会变暗、偏移甚至找不到；如果狭缝过窄，谱线虽细但亮度不足；如果临界反向位置判断不准，折射率结果会直接发生系统偏移。因此我认为这部分的关键操作是“先让谱线稳定清晰，再慢慢寻找反向临界点”，而不是急于读数。

布儒斯特角法让我认识到，操作评分不只看是否写出 \(n=\tan\theta_B\)，也看是否理解“反射光最暗”的物理含义。最暗位置并不等于屏幕上完全没有光，而是 P 分量反射最弱的位置附近，需要在转台角度和偏振片角度之间反复微调。这个过程比最小偏向角法更依赖观察者判断，所以测得的折射率与最小偏向角法存在差异是可以理解的。

光谱仪和特斯拉线圈部分则体现了实验中的另一类能力：知道仪器显示背后的物理条件。光谱仪不是只要连上电脑就能得到可靠数据，还必须正确调曝光、选定校准峰并判断峰位；特斯拉线圈虽然现象直观，但需要明确高频高压、电磁感应和气体电离的安全边界。综合来看，本实验把几何光学、偏振、光谱测量和高压放电串联起来，要求在课堂中既会操作，也能解释每一步为什么这样做。

\section{可提交性说明}

如果平台“操作评分窗口”只允许上传一个文件，本 PDF 已包含操作完成证据、关键原始数据、数据处理和课后总结，适合作为主提交文件。若平台允许多个附件，可同时附上课堂原始照片或原始数据截图，作为对第4节数据的补充证明；但不建议把完整实验报告直接上传到该窗口，以免与“报告提交窗口”混淆。

\end{document}
